\section{相关工作}

\subsection{全景图像语义分割}

全景图像分割是计算机视觉中的重要任务，在自动驾驶、机器人导航和虚拟现实等领域有广泛应用\cite{ma2021densepass,yang2021omnirange}。全景图像通常采用等距柱状投影(equirectangular projection)表示360°场景，这种投影方式在两极区域会产生严重的几何畸变，给传统卷积神经网络带来挑战\cite{su2017learning}。

早期方法通过设计专门的网络结构来处理全景图畸变。SphereNet提出球面卷积核，直接在球面坐标系中进行卷积操作\cite{su2017learning}。可变形卷积网络(DCN)\cite{dai2017deformable}通过学习自适应的感受野来适应畸变特征，在多种视觉任务中展现了强大的建模能力。基于全景图的特殊性，一些工作设计了专门的投影变换和核变换方法来显式建模几何畸变。

近年来，基于Transformer的方法在全景分割中取得了显著进展。Vision Transformer (ViT)\cite{dosovitskiy2020image}首次证明了纯Transformer架构在图像识别中的有效性，启发了后续大量工作。Swin Transformer\cite{liu2021swin}通过分层架构和移位窗口机制实现了线性复杂度，在多项视觉任务中取得突破。Trans4PASS\cite{zhang2022trans4pass}首次将Transformer引入全景分割，利用自注意力机制建模长距离依赖关系，有效捕获全景图的全局上下文。PanoSwin\cite{yang2023panoswin}在Swin Transformer的基础上设计了全景特定的窗口注意力，通过循环移位窗口处理全景图的周期性边界。Pyramid Vision Transformer (PVT)\cite{wang2021pyramid}提出了金字塔结构的视觉Transformer，为密集预测任务提供了多尺度特征。SegFormer\cite{xie2021segformer}结合Transformer和轻量级MLP解码器，在语义分割中取得了优异的效率-精度权衡。Segmenter\cite{strudel2021segmenter}直接将ViT应用于语义分割，验证了Transformer在密集预测任务中的潜力。CrossViT\cite{chen2021crossvit}通过交叉注意力融合多尺度特征，进一步提升了分割性能。

数据集方面，DensePASS\cite{ma2021densepass}提供了大规模的密集标注全景分割数据集，包含19,657张街景全景图像，推动了该领域的快速发展。Stanford2D3D和Matterport3D也为全景场景理解提供了丰富的训练数据。Cityscapes\cite{cordts2016cityscapes}作为透视图像分割的经典数据集，常被用作全景域适应的源域。

\subsection{无监督域适应}

无监督域适应(Unsupervised Domain Adaptation, UDA)旨在将模型从有标注的源域迁移到无标注的目标域，减少标注成本\cite{tsai2018learning,zou2018unsupervised}。语义分割的域适应尤其具有挑战性，因为需要在像素级别对齐不同域之间的分布差异。

\textbf{对抗训练方法}通过域判别器使特征分布在域之间不可区分。AdaptSegNet\cite{tsai2018adaptsegnet}在输出空间进行对抗对齐，开创了语义分割域适应的先河。ADVENT\cite{vu2019advent}通过熵最小化进一步提升对齐效果，利用目标域的熵分布作为自监督信号。这些方法通过对抗训练缩小域间差异，但训练过程往往不稳定，需要精心调整超参数。

\textbf{自训练方法}使用模型预测的伪标签进行训练，逐步提升目标域性能。CBST\cite{zou2018cbst}通过类别平衡的伪标签选择避免类别偏差，确保稀有类别也能得到充分训练。该方法简单有效，成为域适应的重要基线。原型对齐方法\cite{zhang2021prototypical}通过对齐源域和目标域的类别原型减少域偏移，ProDA利用原型引导的伪标签细化提升了目标域的分割质量。

\textbf{风格迁移方法}将源域图像转换到目标域风格，从而缩小视觉外观差异。FDA通过傅里叶域适应实现高效的风格迁移，避免了复杂的生成对抗网络训练。这类方法在保持语义内容的同时改变图像风格，为域适应提供了新的视角。

\textbf{混合增强方法}通过源域和目标域的图像混合生成高质量训练样本。DACS\cite{tranheden2021dacs}通过跨域混合采样生成伪标签，有效利用了源域和目标域的互补信息。DAFormer\cite{hoyer2022daformer}结合Transformer架构和稀有类采样进一步提升性能，在多个域适应基准上取得了最先进的结果。HRDA\cite{hoyer2022hrda}通过上下文感知的高分辨率域适应处理多尺度信息，在保持高分辨率的同时实现了高效训练。MIC\cite{hoyer2023mic}利用掩码图像一致性增强上下文建模，进一步推动了域适应的性能边界。

针对全景域适应的特殊挑战，一些工作设计了专门的方法。DPPass\cite{zheng2023dppass}同时考虑风格和畸变差异，通过双路径域适应处理全景图的特殊性。DATR\cite{zheng2023datr}利用邻域信息进行畸变感知的域适应。GoodSAM\cite{zhang2024goodsam}通过Segment Anything Model桥接域和容量差异，在全景分割中取得了优异性能。SFUDA360\cite{zheng2024sfuda360}和SFUDA360++\cite{zheng2025sfuda360pp}探索了无源域的全景分割域适应，通过学习可靠的类别原型实现了无需源域数据的适应。

最近，基于大模型的方法开始出现。OmniSAM\cite{zhong2025omnisam}首次将SAM2的记忆机制引入全景域适应，通过滑动窗口和记忆注意力处理高分辨率全景图。然而，其训练成本较高，限制了实际应用，这正是本文重点解决的问题。

\subsection{视觉大模型与记忆机制}

视觉大模型通过大规模预训练在多种视觉任务上展现了强大的泛化能力。Mask R-CNN\cite{he2017mask}首次将实例分割与目标检测统一，为后续分割模型奠定了基础。Mask2Former\cite{cheng2022masked}通过掩码注意力机制统一了全景、实例和语义分割，展现了Transformer在分割任务中的潜力。DeepLab系列\cite{chen2017deeplab}通过空洞卷积和全连接CRF实现了高效的语义分割，成为该领域的经典方法。PSPNet\cite{zhao2017pyramid}通过金字塔场景解析网络捕获多尺度上下文信息，在多个分割基准上取得了优异性能。

Segment Anything Model (SAM)\cite{kirillov2023segment}通过在SA-1B数据集(包含1100万张图像和11亿个掩码)上预训练，实现了零样本分割能力。SAM采用提示编码器-图像编码器-掩码解码器的架构，可以通过点、框或文本提示进行交互式分割。这一突破性工作展示了大规模预训练在分割任务中的巨大潜力。

SAM2\cite{ravi2025sam2}将SAM扩展到视频域，引入了记忆编码器和记忆注意力机制。记忆机制的核心思想是存储历史帧的特征表示，通过交叉注意力机制使当前帧查询能够访问历史信息，从而提升时序一致性。然而，随着序列长度增加，记忆Token规模线性增长，导致注意力计算成本急剧上升。对于长序列视频或多窗口全景图，这一问题尤为严重。

记忆机制在其他视觉任务中也得到了广泛应用。Video Object Segmentation中的方法通过记忆网络跟踪目标，实现了高精度的视频分割。Tracking领域也采用了类似的记忆-查询机制来处理长时序信息。然而，这些方法大多未专门优化记忆的计算效率，在处理高分辨率或长序列时面临性能瓶颈。

\subsection{高效注意力机制}

降低Transformer计算复杂度是深度学习领域的重要研究方向。标准自注意力\cite{vaswani2017attention}的计算复杂度为$O(N^2)$，其中$N$为Token数量，在高分辨率视觉任务中成本高昂。Transformer架构\cite{vaswani2017attention}最初在自然语言处理中取得成功，后被广泛应用于计算机视觉任务。

\textbf{稀疏注意力方法}通过限制注意力范围降低复杂度。Sparse Transformer\cite{child2019generating}使用固定的稀疏模式，只关注部分位置的Token，将复杂度降至$O(N\sqrt{N})$。Swin Transformer\cite{liu2021swin}通过窗口注意力和移位窗口机制实现线性复杂度，在保持性能的同时大幅降低计算量。这类方法通过限制注意力的空间范围实现加速，但可能损失全局建模能力。

\textbf{线性注意力方法}将注意力机制重新表述为核函数形式，实现$O(N)$复杂度。Linear Transformer\cite{katharopoulos2020transformers}通过核技巧将softmax注意力转化为线性操作，避免了显式计算注意力矩阵。这类方法在理论上更加高效，但在实践中可能需要较大的隐藏维度才能达到与标准注意力相当的性能。

\textbf{Token剪枝方法}通过动态移除冗余Token降低计算量。DynamicViT\cite{rao2021dynamicvit}预测Token重要性并动态剪枝，在保持精度的同时减少了计算量。该方法通过轻量级的重要性预测器判断哪些Token可以被安全移除。EViT保留类别Token并将剪枝Token的信息融合到保留Token中，确保信息不丢失。Token Merging (ToMe)\cite{bolya2023tome}通过合并相似Token减少序列长度，这一简单有效的策略在ViT中取得了显著的加速效果。Evo-ViT提出渐进式Token剪枝和慢速-快速更新策略，逐层减少Token数量。REST\cite{zhang2021rest}通过重用和共享Token实现高效的视觉识别。Token-to-Token ViT (T2T-ViT)\cite{yuan2021tokens}通过渐进式Token化过程建模局部结构，提升了ViT的效率和性能。内容感知Token共享\cite{lu2023tokensharing}专门针对语义分割设计了Token共享策略，在保持分割质量的同时减少了冗余计算。

\textbf{混合方法}结合多种技术实现更好的效率-精度权衡。Multiscale Vision Transformers (MViT)\cite{fan2021multiscale}通过多尺度特征金字塔和池化注意力实现了高效的视频理解。这些方法综合利用了稀疏注意力、Token剪枝和架构优化，在多个任务中取得了优异的效率-精度平衡。

本文借鉴Token剪枝和稀疏注意力的思想，针对记忆注意力机制设计了专门的优化方案：通过空间压缩（MTC）降低单个记忆的Token数量，通过查询引导的稀疏检索（SMR）减少参与注意力计算的记忆数量，在保持精度的同时大幅提升训练和推理效率。与现有工作不同，我们的方法专门针对全景域适应中的记忆机制设计，同时提出了循环记忆（SCM）来处理全景图的边界连续性问题，实现了效率与精度的双重优化。
